\chapter{Dos velocidades}

\label{cap:conclusion}
\section{Lo que ocurrió}

El Boletín Oficial de Salta publicó su primer número el 15 de octubre de 1908. En septiembre de 2026, la Provincia dio de baja un contrato de defensas sobre el río La Caldera por doscientos dieciséis millones de pesos \textbf{porque el cauce se había corrido}, la misma razón que en 1940 había obligado a suspender otras obras. Entre una fecha y la otra hay ciento diecisiete años, y el río es el mismo.

Lo que cambió en el medio, y sobre todo en el último cuarto de siglo, es lo que sigue.

El departamento La Caldera fue, entre 2010 y 2022, el cuarto que más creció en población de toda la Argentina: 58,4 por ciento, hasta 12.299 habitantes. Sólo lo superaron dos partidos del frente de urbanización bonaerense y Añelo, la capital de Vaca Muerta ---el ordenamiento se verificó sobre los resultados definitivos de veintitrés de las veinticuatro jurisdicciones; el capítulo \ref{cap:poblacion} precisa el alcance de ese cotejo---. Creció más del triple que el departamento Capital. Lo administran dos municipios clasificados en tercera categoría.

Ese municipio ---al que su propio plan atribuye mil doscientos kilómetros cuadrados, más que los mil noventa y dos que el INDEC asigna a todo el departamento (capítulo \ref{cap:fincas})--- recibió en 2015 un plan urbano ambiental, financiado por el Consejo Federal de Inversiones y asistido técnicamente por la Provincia, cuyo informe final nombra en su primer párrafo a un municipio distinto, situado a cuatrocientos kilómetros; al año siguiente lo reglamentó con un Código de ciento cincuenta y dos artículos. Ese Código clasifica el suelo con precisión, declara no urbanizable la ribera del río, nombra uno por uno los barrios que considera espontáneos, ordena gravar de manera permanente los loteos en áreas no aptas y somete la venta de parcelas a garantía de obra y licencia municipal previa. El plan había dejado como proyecto pendiente el estudio que determina hasta qué cota es seguro construir.

En 2020, cinco años después de recibir el plan, el mismo municipio declaró públicamente que la coparticipación no le alcanzaba para pagar los sueldos de su personal.

Entre una cosa y la otra ocurrió todo lo demás. Los loteos que el Código llama no planificados se inundaron dos veces y una sentencia declaró la obligación estatal de proteger y recomponer la cuenca, con criterio ecocéntrico y cita de la Corte Interamericana. La orden se reabsorbió y reapareció cuatro años más tarde como propuesta de gestión del propio obligado. El Plan se aprobó en febrero de 2026, con el municipio intimado bajo apercibimiento de remisión a la justicia penal, más de siete años después de los hechos, en una causa que la Constitución define como vía expedita. El desarrollador de uno de esos loteos quedó identificado en un boletín oficial y figura como tercero en el expediente. Un loteo cuya urbanización el propio municipio declaró incompatible con su Código recibió, en el mismo acto, la indicación del procedimiento para urbanizarlo igual. Un tercer loteo, cerrado, cedió su ingreso sobre una ruta nacional.

Al final del período, el departamento tenía una de cada cinco viviendas sin nadie viviendo en ella, más del doble de la proporción provincial. Al comienzo también: el ciclo duplicó el parque habitacional sin alterar esa proporción. En esos mismos años perdió el 88 por ciento de su superficie agropecuaria declarada ---en su mayor parte monte y pastizal, según el censo de 2008---, mientras la provincia ganaba.

Simultáneamente, bajo la misma autoridad provincial, funcionó el otro circuito. En un barrio de la capital contiguo al departamento, donde según su centro vecinal buena parte de los adjudicatarios son caldereños, familias que ocupaban terrenos fiscales obtuvieron en 2025 su adjudicación en venta ---en un caso, tras un expediente iniciado en 1995---, en cien cuotas y con hipoteca, después de que Inmuebles y el Instituto Provincial de Vivienda informaran que no tenían otro inmueble ni otra solución habitacional, y con la obligación de habitar de manera continua bajo pena de perder, sin indemnización, el lote y todo lo construido en él. Y quien no ocupa y pide un lote fiscal por sorteo debe acreditar residencia policial, declarar bajo pena de falso testimonio, dar referencias vecinales con nombre y apellido y una franja horaria de visita domiciliaria, ser argentino nativo y no tener inmuebles en ningún lugar del mundo.

\section{Las tres tesis, al cabo}

\textbf{Los dos circuitos no se distinguen por su legalidad.} Se distinguen por la intensidad con que el Estado se ocupa de cada uno, y esa intensidad es inversamente proporcional a la capacidad de sus destinatarios. Al que menos tiene se le exige más, se lo controla mejor y se lo hace esperar décadas. Al que más puede se le exige un procedimiento que no se verifica, se lo controla mediante un aparato que no existe y se lo atiende en el plazo de un ciclo comercial.

\textbf{Y la opacidad que sostiene esa asimetría no está en el acto sino en el criterio.} Seis instrumentos de dos niveles de gobierno publican la envoltura y reservan a un anexo el contenido que permitiría verificar si alguien fue tratado con la misma vara que otro. La zonificación que prohíbe prevé el trámite que autoriza. Y los actos que el Código exige antes de vender una parcela ---garantía de obra y licencia de comercialización, entre otros, a los que el decreto provincial de 2019 suma una autorización de preventa con caución--- no constan en ninguna fuente pública para ninguno de los emprendimientos relevados.

\textbf{Y ninguna de las dos cosas es reciente.} Las unidades sobre las que hoy se lotea son los partidos policiales de 1909. La forma institucional del municipio —comisión con presidente designado, presupuesto aprobado desde afuera— es la de 1908 y la de 1922. La falta de mediciones fue declarada por escrito por el propio Estado en 1924, y las series que se levantaron después ---entre 1942 y 1986--- nunca alcanzaron el umbral que exigiría determinar hoy las líneas de ribera por aforo. Y el pleito que en 1908 se resolvió porque el título no aparecía prefigura, con exactitud incómoda, la razón por la cual este libro no puede establecer si una matrícula proviene de otra.

El mercado inmobiliario no construyó este dispositivo. Lo encontró armado.

\begin{cautela}
Este resultado no requiere que nadie haya decidido producirlo, y ése es su rasgo más incómodo. No hay en el material relevado prueba de captura, de connivencia ni de blanqueo. Hay una Provincia que transfirió competencias sin transferir capacidad; un municipio que aceptó una norma diseñada para un aparato que no tiene; un mercado que ocupó el espacio que esa incapacidad dejó libre; y un régimen de tierra fiscal minucioso, escrito con buenas intenciones y aplicado sobre la población que menos margen tiene para incumplirlo.

La responsabilidad política existe y es identificable, pero no está donde la denuncia la buscaría. Está en la decisión sostenida de asignar al municipio la función de regular el suelo del área metropolitana sin darle los medios de hacerlo, y en la de no publicar junto a las normas los anexos que las vuelven verificables.

Y este libro no puede distinguir, con los documentos que consiguió, entre un Estado desbordado y un Estado desbordado que a alguien le conviene que siga desbordado.
\end{cautela}

\subsection*{Cuánto hay que oscurecer el cuadro}

Entre «el Estado falla» y «alguien lo arregló» hay una escala, y conviene decir en qué punto de
esa escala deja el material, porque el lector va a hacer esa cuenta igual.

\textbf{Antes de la escala hay que separar dos cosas que el lector va a juntar.} Este libro
mide dos cuadros distintos y no se mueven juntos. Uno es el de las consecuencias, y está
establecido con cifras que el apartado siguiente detalla: vivienda vacía que más que duplica la
tasa provincial, superficie agropecuaria perdida mientras la provincia ganaba, y una parte
grande de la población sin agua de red en un departamento cuyos ríos se concesionan. Son datos
censales y agropecuarios, publicados con su fuente y su procesamiento a la vista, y
\textbf{no dependen de que nadie haya querido producirlos}. El otro cuadro es el de la
intención, y está indeterminado: no apareció un acto ilegal, ni un pago, ni una instrucción, ni
un acto dirigido a beneficiar a alguien identificable.

\textbf{Un cuadro pesado de consecuencias no prueba una intención, y una intención no probada
no aligera las consecuencias.} Las dos afirmaciones tienen distinto estatuto y conviene no
canjearlas. La primera se sostiene en mediciones que cualquiera puede rehacer con los mismos
datos. La segunda se sostiene en una ausencia, y una ausencia sólo vale dentro de la ventana
que efectivamente se leyó (apéndice \ref{cap:fuentes}): decir que no apareció no es lo mismo
que decir que no hay. Leer la primera como demostración de la segunda es el error que este
capítulo pide no cometer; leer la segunda como atenuante de la primera es el error simétrico, y
es el más cómodo de los dos.

\textbf{Lo que sigue mide el segundo cuadro.} El primero no está en escala: está medido.

\textbf{No está en el extremo claro.} Cuatro cosas lo impiden. La opacidad no está en cualquier
lugar sino exactamente donde deja de ser controlable: se publica el marco y se reserva el
criterio. Las líneas de ribera se publican dos días y no se asientan en la matrícula, de modo
que el que compra no puede enterarse y el que vende sí sabe. La asimetría entre los dos
circuitos es constante y va siempre en la misma dirección. Y la concentración de funciones no
empieza con el mercado: en 1909 el que clasificaba el impuesto del departamento mandaba su
policía, el juez de paz suplente era a la vez comisario suplente, el comisario de un partido era
propietario en ese partido y el presidente municipal administraba la finca más grande del
departamento (capítulos \ref{cap:fincas} y \ref{cap:siglo}).

\textbf{Y no está en el extremo oscuro.} Tres cosas lo impiden, y son del mismo peso. No
apareció un solo acto ilegal. No apareció un pago, una instrucción ni un acto dirigido a
beneficiar a alguien identificable. Y donde parecía haber una decisión deliberada, el archivo
fue en contra: la negativa al puente de 1935 cae dentro de un ajuste presupuestario de Vialidad
que la explica sin otra razón, y en 1940 el propio Estado puso ese puente en su plan de quince
años (capítulos \ref{cap:defensas} y \ref{cap:infraestructura}).

\textbf{Qué movería el cuadro, y es verificable.} Tres documentos que están pedidos y no
llegaron: las fojas del expediente de El Durazno, los antecedentes dominiales de las matrículas
del corredor y las respuestas a las cartas de réplica. Si en alguno de ellos aparece una
instrucción, una oposición registrada de propietarios a una obra pública o una partida desviada
a pedido de alguien, \textbf{el cuadro se oscurece y, al mismo tiempo, el término del título deja de
corresponder}: habría estratega, y entonces lo que hay que explicar no es un dispositivo sino
una decisión. Las tres tesis tienen sus propios refutadores (capítulo \ref{cap:planteo}); lo que caería es la palabra que las reúne, y el capítulo \ref{cap:prospectiva} enuncia esa condición de refutación en sus
propios términos. Hasta hoy no apareció.

\section{Lo que este libro no probó}

De las cuatro hipótesis heredadas, H.20 se confirma sólo en su término de segregación residencial; H.21 no fue investigada; H.22 tiene una verificación parcial en un municipio vecino; H.23 tiene un indicio periodístico sin respaldo documental. La hipótesis del blanqueo, que motivó el proyecto, no encontró un solo elemento de prueba.

Se estableció que el departamento tiene 4.991 viviendas particulares y 4.003 hogares --- 988 sin residentes, el 19,80 por ciento del parque, contra el 8,55 provincial ---, que una de sus tres fracciones censales llega al 43,17 por ciento, y que esa tasa era del 19,01 por ciento ya en 2010: la vivienda vacía es anterior al ciclo de loteos y no fue producida por él. No se estableció qué proporción de esas viviendas son segundas residencias, cuáles están en venta y cuáles en construcción.

Tampoco se probó que el departamento fuera lugar de veraneo de la elite salteña. Dos vecinos lo afirman y agregan un mecanismo: sin puente, las crecidas del verano dejaban a los veraneantes aislados de la capital. El archivo confirma la premisa material ---en 1935 el pueblo pedía un puente que no tenía, y el río crecía todos los años en la época de lluvias---, pero no el veraneo ni el aislamiento como motivo; y dos documentos contradicen que la falta de puente fuera buscada (capítulo \ref{cap:poblacion}).

No se probó que exista relación causal entre los loteos y los aluviones, aunque sí que la advertencia vecinal sobre el desmonte ribereño fue formulada dos años antes del daño ---y que las advertencias técnicas, del INTA y de organizaciones especializadas, eran anteriores--- y que el estudio de cota máxima edificable figuraba como pendiente en el plan vigente en ese momento.

No se probó la copropiedad del loteo Solanas por parte de la actual titular del organismo de tierra fiscal, y por eso el libro no la afirma. Se probó, en cambio, que hay dos comunidades kollas registradas en el departamento y que la zona alta de la cuenca es territorio relevado de la Kondorwaira; y el cruce catastral mostró que ninguno de los loteos documentados toca la parcela que la comunidad reclama ---el catastro 102, fiscal---, aunque el perímetro exacto sigue en la carpeta técnica. Y el Censo Agropecuario registra que el 46,6 por ciento de las explotaciones carece de límites definidos, sin que este libro haya podido identificar a esa población.

Se estableció, en cambio, que entre 2002 y 2018 el departamento perdió el 88,4 por ciento de su superficie agropecuaria --- 87.357 hectáreas --- mientras la provincia ganaba un 2,8 por ciento, y que la pérdida se concentró en las explotaciones mayores a mil hectáreas, de las que sólo sobrevive una. No se estableció qué ocurrió con esa tierra. Se estableció, sí, qué se dejó de sembrar: las legumbres desaparecieron por completo, las forrajeras cayeron a la mitad y las hortalizas más del cincuenta por ciento, mientras los cultivos industriales crecían un 11,7 por ciento. Desapareció la agricultura mixta, no la comercial.

Quedó refutada, en cambio, la hipótesis de que la escasez de leña hubiera provocado el colapso de la actividad tabacalera en el período intercensal: la superficie de cultivos industriales creció. Y quedó además acotada por otra vía: en las dos puntas del período, el departamento aparece en la cola del mapa tabacalero provincial, de modo que un cultivo que nunca fue dominante allí difícilmente pudo arrastrar consigo el 88 por ciento de la superficie.

\section{Una consecuencia presente}

Un dato del Anexo I de la Ley 8483, cuyo marco analiza el capítulo \ref{cap:opacidad}, no describe el pasado sino que rige hoy, y por eso se repite aquí.

La revisión del Ordenamiento Territorial de Bosques Nativos aprobada por la Ley provincial 8483, publicada el 3 de enero de 2025, asigna a la cuenca Mojotoro--Lavayén --- que contiene la microcuenca del río La Caldera --- veintinueve mil setecientas cuarenta y cuatro hectáreas de Área de Producción y Conservación y \textbf{cero hectáreas de verde}. Cero coma cero por ciento categorizable.

Conforme el criterio establecido en el propio informe técnico de la ley, las cuencas excedidas en su umbral máximo de deforestación no son categorizables para proyectos de cambio de uso del suelo, y en ellas la deforestación queda prohibida.

La Tabla 17 del informe marca esa cuenca como excedida, y la adenda de la autoridad provincial confirma el cero por ciento. El desmonte de bosque nativo está prohibido en ella desde enero de 2025 por aplicación del ordenamiento provincial. Según las capas del propio informe, esa cuenca cubre tres cuartos del departamento y contiene entre el 98 y el 99 por ciento de su Área de Producción y Conservación: en la práctica, todo el bosque nativo que el ordenamiento reconoce en La Caldera está bajo un régimen que prohíbe el desmonte (capítulo \ref{cap:opacidad}). Entre una quinta y una cuarta parte del departamento queda sin categoría, porque el mapa no la cubre.

\section{Lo que falta, y no se consigue buscando}

Faltan cinco de los seis contenidos reservados que el capítulo \ref{cap:opacidad} inventaría, el informe de dominio de las matrículas 4.366 y 6.558, el texto íntegro de la sentencia, las intendencias anteriores a 1999 --- que para 1983--1987 no se buscan como proclamaciones sino como decretos de designación, porque entonces los intendentes no se elegían --- y el perímetro fino del territorio comunitario relevado, que el cruce a nivel de catastro no reemplaza.

Y falta lo que no está en ningún archivo. Este libro no contiene entrevistas propias: sólo dos conversaciones informales con vecinos sobre el veraneo, registradas como testimonio. El amparista que sostuvo siete años de litigio, la bióloga que documentó el desmonte, las profesionales del COPAIPA que actuaron como veedoras, el centro vecinal de La Calderilla, las familias del barrio Juan Manuel de Rosas que esperaron años ---en un caso, treinta--- su adjudicación: todos son localizables y ninguno habló con este trabajo. Tampoco lo hicieron los desarrolladores, cuya versión aparece únicamente en las presentaciones de sus abogados dentro de notas periodísticas. Ni las familias de la Comunidad Kolla Kondorwaira ---veintisiete según el parte provincial de 2017, treinta y una según el artículo de 2024---, que habitan la cabecera de esta cuenca y que \textbf{escribieron su propio libro sobre ella} --- de modo que en su caso la voz no falta: falta que este trabajo haya ido a buscarla.

\begin{cautela}
Un relevamiento documental establece qué dicen las normas y los expedientes. No establece qué ocurre cuando una familia entra a una oficina municipal, ni qué se le responde, ni cuánto tarda, ni quién la atiende. Y el capítulo \ref{cap:opacidad} mostró que ése es precisamente el lugar donde el criterio se aplica sin registrarse: la oficina donde la excepción se tramita y el anexo se consulta.

Es decir que las dos carencias de este libro --- la falta de entrevistas y la imposibilidad de distinguir el desborde del desborde conveniente --- son la misma carencia. La segunda sólo se resuelve por la vía de la primera.

Ése es el trabajo que sigue.
\end{cautela}
